% Chapter 4 -- Claim: A Profitable Backtest Means a Profitable Strategy
% Requires: amsmath, amssymb, booktabs, array, tabularx, graphicx, tcolorbox (breakable, skins), hyperref or url.
% Defines the callout box only if the preamble has not.
\providecolor{fenavy}{HTML}{1B2A6B}
\providecolor{feshade}{HTML}{F3F5FA}
\providecolor{feorange}{HTML}{B8481C}
\makeatletter
\@ifundefined{callout}{%
  \newtcolorbox{callout}{breakable, enhanced,
    colback=white, colframe=gray!70, boxrule=0.5pt, arc=1mm,
    borderline west={2.5pt}{0pt}{gray!70},
    left=8pt, right=8pt, top=4pt, bottom=4pt}%
}{}
\makeatother

\chapter{Claim: A Profitable Backtest Means a Profitable Strategy}
\label{ch:04}

Chapter~\ref{ch:03} described the factor zoo: hundreds of published patterns in average returns that the CAPM cannot explain. Almost all of them were documented with \textbf{paper returns}\index{return!paper versus implementable}, the returns of portfolios that are rebuilt every month at recorded closing prices, in any size, at no cost. Researchers have long known that this is a simplification, and the reason it matters goes back to the definition of an efficient market. Jensen (1978) argued that a market is efficient with respect to some information if no one can earn \emph{economic} profits by trading on it, which means profits after all the costs of trading. Grossman and Stiglitz (1980) went further. They showed that prices cannot reflect all information perfectly, because then no one would be paid for the cost of gathering it. On this view, some return patterns should survive on paper but not after costs, since costs are exactly what keeps them from being traded away.

The question became pressing in the early 2000s, when researchers began to price the trades behind the best-known anomalies\index{anomaly}. Lesmond, Schill and Zhou (2004) found that momentum\index{momentum} trades concentrate in stocks that are expensive to trade, and concluded that its profits largely disappear after costs. Korajczyk and Sadka (2004) measured how prices move when large orders are executed. They found that momentum profits shrink as the size of the portfolio grows, and that even versions designed to trade cheaply could absorb only about \$5 billion, at December 1999 market values, before the profits vanished. Novy-Marx and Velikov (2016) then estimated trading costs for 23 well-known anomalies and found that most strategies that trade more than half of their portfolio each month lose their significance after costs. Frazzini, Israel and Moskowitz (2018) challenged the size of these estimates. Using a large manager's own trading records, they found costs far below what the academic estimates assumed. The same patterns therefore look profitable or unprofitable depending on whose costs are used.

The answer matters for both research and practice. For research, it decides whether an anomaly is evidence against market efficiency or only a pattern that no investor could exploit. For practice, it decides which published strategies are worth running, how much money each can manage, and whether a fund that offers them can deliver its backtest. The claim on trial\index{claims-on-trial method} is the assumption behind most published backtests\index{backtest}: \textbf{a strategy that is profitable in a backtest is profitable to trade.}

The chapter draws on two papers about the cost of trading. Novy-Marx and Velikov (2016) estimate how much of each published anomaly survives trading costs. Frazzini, Israel and Moskowitz (2018) use real institutional trades to show that costs can be lower than academic estimates, but that they grow with the size of the fund. The lab in Section~\ref{sec:04-lab} replicates the main calculations of both papers in order to find out which strategies stay profitable after costs, what cost level would wipe out each one, and how much money each can manage.

By the end of the chapter you will be able to:

\begin{itemize}
  \item read a limit order book\index{order book}, and explain bid--ask spreads\index{bid--ask spread}, market orders\index{market order} and limit orders\index{limit order};
  \item explain how long and short positions work, what short selling\index{short selling} costs, and what a zero-investment portfolio\index{long--short portfolio} is;
  \item break trading costs into commissions, bid--ask spreads, market impact\index{market impact} and the costs of short selling, and estimate the main ones from public data;
  \item compute a strategy's turnover\index{turnover} and explain what drives it;
  \item convert gross returns into net returns, and compute the break-even cost\index{transaction costs!break-even} a strategy can bear;
  \item estimate how market impact limits the money a strategy can manage;
  \item design strategies that trade less without giving up much of their signal.
\end{itemize}

In the lab you will take three of the best-known long--short strategies, value, momentum and short-term reversal\index{short-term reversal}, measure how much each trades, charge each for trading, and see how their ranking changes. A simulated market with known expected returns\index{return!expected} then shows how much a cost-aware design can recover.

\section{Foundations: how stocks are traded and sold short}
\label{sec:04-foundations}

\subsection{The order book}

Most stocks trade on electronic exchanges. Each exchange keeps a \textbf{limit order book}\index{order book|textbf}, which is a list of standing offers to buy and to sell. The offers to buy are called \textbf{bids}, and the offers to sell are called \textbf{asks}. Each offer states a price and a number of shares (Figure~\ref{fig:04-book}).

The highest bid and the lowest ask are the \textbf{best quotes}. The gap between them is the \textbf{bid--ask spread}\index{bid--ask spread|textbf}, and the price halfway between them is the \textbf{midpoint}\index{midpoint|textbf}. Spreads are usually quoted in \textbf{basis points}\index{basis point|textbf} of the midpoint. A basis point is a hundredth of a percent, so 100 basis points make 1\%.

\begin{figure}[htbp]
  \centering
  \scalebox{0.8}{%
\begin{tikzpicture}[
    x=0.014cm, y=0.5cm,     % x: shares; y: price levels
    ask/.style={fill=feorange!35},
    bid/.style={fill=fenavy!75},
    price/.style={anchor=east, font=\small},
  ]
    % Asks (sellers) above the midpoint: level / price / shares
    \foreach \wfL/\wfP/\wfQ in {1/50.02/200, 2/50.03/350, 3/50.04/300,
                                4/50.05/450, 5/50.06/500} {
      \fill[ask] (0,\wfL+0.5-0.38) rectangle (\wfQ,\wfL+0.5+0.38);
      \node[price] at (-10,\wfL+0.5) {\$\wfP};
    }
    % Bids (buyers) below the midpoint
    \foreach \wfL/\wfP/\wfQ in {1/49.98/250, 2/49.97/300, 3/49.96/400,
                                4/49.95/350, 5/49.94/550} {
      \fill[bid] (0,-\wfL-0.5-0.38) rectangle (\wfQ,-\wfL-0.5+0.38);
      \node[price] at (-10,-\wfL-0.5) {\$\wfP};
    }

    % Best ask and best bid
    \draw[feorange, thick] (0,1.12) rectangle (200,1.88);
    \draw[fenavy, thick]   (0,-1.88) rectangle (250,-1.12);
    \node[feorange, anchor=west, font=\small] at (215,1.5)
      {best ask: buyers pay \$50.02};
    \node[fenavy, anchor=west, font=\small] at (265,-1.5)
      {best bid: sellers get \$49.98};

    % Midpoint
    \draw[densely dashed] (-10,0) -- (560,0);
    \node[anchor=south east, font=\small, inner sep=2pt] at (560,0) {midpoint \$50.00};

    % Spread bracket
    \draw[decorate, decoration={brace, amplitude=4pt}]
      (-150,-1.5) -- (-150,1.5)
      node[midway, anchor=east, xshift=-5pt, align=right, font=\small]
      {spread\\\$0.04\\(8 bp)};

    % Headings and size axis
    \draw[thin] (0,-6.1) -- (0,6.1);
    \node[feorange, font=\small\bfseries, anchor=south west] at (0,6.1)
      {Asks: sellers' limit orders};
    \draw[-stealth, thin] (0,-6.1) -- (580,-6.1)
      node[anchor=west, font=\small] {shares};
    \node[fenavy, font=\small\bfseries, anchor=north west] at (0,-6.3)
      {Bids: buyers' limit orders};
  \end{tikzpicture}%
}
  \caption{A limit order book. Sellers' asks sit above the midpoint and buyers'
    bids below it. A buyer who wants to trade at once pays the best ask,
    \$50.02; a seller receives the best bid, \$49.98. The gap of \$0.04,
    8 basis points of the \$50.00 midpoint, is the bid--ask spread.}
  \label{fig:04-book}
\end{figure}

\textbf{Worked example.} In Figure~\ref{fig:04-book}, the best bid is \$49.98 and the best ask is \$50.02. The midpoint is \$50.00 and the spread is \$0.04, which is $0.04/50 = 0.0008$, or 8 basis points. An investor who buys at the ask and sells again at once at the bid loses those 8 basis points. The cost of one side of the trade, called the \textbf{one-way cost}\index{transaction costs!one-way and round-trip}, is half the spread: 4 basis points.

\subsection{Market and limit orders}

There are two basic ways to trade.

\begin{itemize}
  \item A \textbf{market order}\index{market order|textbf} trades at once against the best available quote. A buyer pays the ask and a seller receives the bid. A market order \emph{takes} liquidity\index{liquidity}, and it pays the spread for doing so.
  \item A \textbf{limit order}\index{limit order|textbf} names a price and waits in the book until someone trades against it. A limit order \emph{supplies} liquidity and may earn the spread. Its risk is that it may never be filled, or that it is filled just before the price moves against it.
\end{itemize}

Large orders are usually split into many small ones and executed over hours or days, so that each piece moves the price less.

\subsection{Who provides liquidity}

\textbf{Market makers}\index{market maker|textbf} and high-frequency traders post bids and asks all day and earn the spread. In return they bear two risks. The first is \textbf{inventory risk}: prices may move while they hold shares they have bought but not yet sold. The second is the risk of trading with investors who know more than they do. Both risks are larger for small, volatile and thinly traded stocks, which is why spreads are wider for those stocks.

\subsection{Long and short positions}

A \textbf{long} position owns a security and gains when its price rises. A \textbf{short} position borrows shares, sells them, and later buys them back to return them. It gains when the price falls.

Short selling has costs that a long position does not. The short seller must pay any dividends to the lender. The short seller must also post \textbf{collateral}\index{collateral|textbf}, cash or securities held as security for the loan, and pay a \textbf{borrowing fee}\index{borrowing fee|textbf} set in the securities-lending market. Most large stocks are cheap to borrow. A minority are hard to borrow and expensive, and the lender can \textbf{recall} the loan at any time, which forces the short seller to buy the shares back.

\subsection{Long--short portfolios}

Academic factors such as HML\index{factor!HML (value)}, introduced in Section~\ref{sec:03-factors}, are \textbf{zero-investment} portfolios. They hold one dollar long and one dollar short, so the money raised by the short sale pays for the long side. A factor's ``return'' is therefore the difference between the returns of the two sides. It is a spread, not a return on money invested. A real long--short fund must hold capital against both sides, and it quotes its returns relative to that capital.

\section{Paper returns versus implementable returns}
\label{sec:04-paper}

The factor returns in the Kenneth French Data Library\index{Kenneth French data library}, like most portfolio returns in academic research, are \textbf{paper returns}\index{return!paper versus implementable|textbf}. Each month the portfolio is rebuilt from the latest sorts, at closing prices, with no charge for trading. This construction is simple and transparent, which is why the data are so useful. But it answers the question ``what pattern exists in returns?'' It does not answer the question ``what could an investor have earned?''

An \textbf{implementable return} is one that an investor could actually have earned. It differs from a paper return in four ways.

\begin{enumerate}
  \item \textbf{Trading costs.}\index{transaction costs|textbf} Every purchase pays more than the midpoint, and every sale receives less. Large orders also move the price.
  \item \textbf{Short-selling costs and limits.} Selling short requires borrowing shares. Borrowing costs a fee, may be impossible for some stocks, and can be recalled.
  \item \textbf{Timing and execution.} A signal computed from closing prices cannot be traded at those same prices. Orders take time to fill, and fast-moving signals lose value while they wait.
  \item \textbf{Scale.} A strategy that works with \$10 million may not work with \$10 billion, because the trades themselves change prices.
\end{enumerate}

None of these makes a paper return wrong. A pattern that cannot be traded profitably may still tell us something about risk or about how prices are formed. But a claim about profits needs an implementable return, and the two can differ by more than the whole premium.

The gap is largest for strategies that trade often, that hold small or illiquid stocks, or that rely on the short side. It is smallest for slow strategies in large stocks. This is one reason the market portfolio is so hard to beat after costs: it hardly trades at all.

\section{The anatomy of trading costs}
\label{sec:04-anatomy}

The cost of a trade has several parts, and each behaves differently.

\begin{center}
\small
\begin{tabularx}{0.9\linewidth}{@{}>{\raggedright\arraybackslash}p{0.19\linewidth}>{\raggedright\arraybackslash}X>{\raggedright\arraybackslash}p{0.32\linewidth}@{}}
\toprule
Component & What it is & How it scales \\
\midrule
Commissions and fees & Payments to brokers and exchanges & Small for institutions today; roughly proportional to the amount traded \\
Bid--ask spread & Buying at the ask and selling at the bid, above and below the midpoint & Proportional to the amount traded; wider for small, volatile and illiquid stocks \\
Market impact & The price moving against an order while it is executed, because the order carries information or demands liquidity & Grows faster than the trade, roughly with the square root of the order's size relative to volume \\
Delay and opportunity cost & Prices moving before the order is filled, or trades not made at all & Largest for fast signals \\
Short-selling costs & Fees to borrow shares, the risk of recall, and rules that restrict shorting & Small for most stocks, very large for a few \\
\bottomrule
\end{tabularx}
\end{center}

The unit used throughout this chapter is the \textbf{one-way cost}\index{transaction costs!one-way and round-trip|textbf}: the cost of buying, or of selling, one dollar of stock, as a fraction of that dollar. For a small trade it is roughly half the bid--ask spread plus commissions. A \textbf{round trip}, buying and later selling, costs twice the one-way cost.

\textbf{Costs have fallen, but unevenly.} In 2001 U.S. stock prices moved from fractions of a dollar to decimals, a change called \textbf{decimalization}\index{decimalization|textbf}. Spreads narrowed sharply after it, and narrowed further as electronic trading spread. Large stocks now trade at spreads of a few basis points, while small, illiquid stocks still cost far more. This matters when we read long histories. A strategy tested from 1963 faced costs in its early decades that were several times today's. Applying today's costs to the whole history flatters the past, and applying the past's costs to today understates what is possible now.

\textbf{Short selling.}\index{short selling|textbf} D'Avolio (2002) studied the market for borrowing stock. He found that the large majority of U.S. stocks were cheap to borrow, but that a small set of ``special'' stocks carried high fees, and that these were often the stocks that short sellers most wanted to sell. Loans can also be recalled at inconvenient times, and regulators can ban short selling outright, as the U.S. Securities and Exchange Commission did for financial stocks in September 2008.

These limits matter for prices as well as for costs. Miller (1977) argued that when pessimists cannot sell short, prices reflect only the views of optimists, so overpricing can last. Stambaugh, Yu and Yuan (2012) found that anomaly returns come more from the short side than from the long side, especially after periods of high investor sentiment, which is where short-sale limits matter most. A long--short paper portfolio that ignores these frictions can overstate what is achievable by much of its premium.

The lab uses a single one-way cost and a single borrowing fee, so that the calculations are easy to follow. Real costs vary across stocks and over time, and Section~\ref{sec:04-measuring} shows how they can be estimated.

\section{Measuring costs from public data}
\label{sec:04-measuring}

The best measure of trading costs is a trader's own record. Perold (1988) proposed the \textbf{implementation shortfall}\index{implementation shortfall|textbf}, which compares the value of a paper portfolio, priced when the decision was made, with the value of the portfolio actually achieved. It captures spreads, impact, delay and trades that were never completed. Most researchers do not have such records, so the field relies on estimates.

\textbf{From quotes and trades.} With intraday data, such as the TAQ (Trade and Quote) database of every U.S. trade and quote, two spreads can be measured.

\begin{itemize}
  \item The \textbf{quoted spread}\index{bid--ask spread!quoted|textbf} is the gap between the best bid and the best ask.
  \item The \textbf{effective spread}\index{bid--ask spread!effective|textbf} is twice the difference between the trade price and the midpoint at the time of the trade.
\end{itemize}

Effective spreads are usually smaller than quoted spreads, because many trades execute between the quotes.

\textbf{From daily data.} Several estimators infer the spread from daily prices alone, which makes them available back to the 1920s.

\begin{itemize}
  \item \textbf{Roll (1984)}\index{bid--ask spread!estimators|textbf} uses \textbf{bid--ask bounce}\index{bid--ask bounce|textbf}. Trades alternate at random between the bid and the ask, so successive price changes are negatively correlated, and the size of that correlation\index{correlation} reveals the spread.
  \item \textbf{Hasbrouck (2009)} estimates a similar model with a Gibbs sampler, a Bayesian simulation method. It works in cases where the Roll estimator fails.
  \item \textbf{Corwin and Schultz (2012)} use daily high and low prices. The high is usually a trade at the ask and the low a trade at the bid. By comparing one-day and two-day ranges, they separate the part of the range that comes from the spread from the part that comes from volatility\index{volatility}.
  \item \textbf{Abdi and Ranaldo (2017)} combine closing, high and low prices to obtain a more precise estimator.
  \item \textbf{Amihud (2002)}\index{illiquidity|textbf} measures \textbf{illiquidity} as the average absolute return per dollar of trading volume. It is a measure of price impact rather than of the spread.
\end{itemize}

These estimators are noisy for any single stock in any single month. Averaged across a portfolio they are useful, and this is how Novy-Marx and Velikov (2016) used them to estimate the costs of their 23 anomalies.

\textbf{Estimates from practice.} Frazzini, Israel and Moskowitz (2018) studied \$1.7 trillion of trades executed by one large institutional manager in 21 developed markets from 1998 to 2016. They found trading costs about an order of magnitude smaller than earlier academic studies suggested, with price impact as the largest component. Their result shows that costs depend on who is trading and how. A patient institution that often supplies liquidity pays far less than the generic investor assumed by spread estimators. The result also comes from one manager's experience, under that manager's own trading methods, so it may not apply to other investors.

A cost assumption is therefore a judgment, not a fact. A careful study reports results across a range of costs. The lab does so, from zero to 100 basis points one way.

\section{Turnover, net returns and break-even costs}
\label{sec:04-turnover}

\textbf{Turnover}\index{turnover|textbf} measures how much of a portfolio is traded in a period. One-sided monthly turnover is the share of the portfolio that is sold, and replaced, each month. It equals half the sum of the absolute changes in portfolio weights, measured against the weights the portfolio would have drifted to without trading. A turnover of 10\% means that a tenth of the portfolio is replaced each month. A value-weighted\index{index weighting!value-weighted} portfolio whose membership never changes has zero turnover, however much its prices move.

Turnover multiplies every per-trade cost, so it is the first number to ask about any strategy. What drives it is how fast the signal changes compared with how often the portfolio is rebuilt.

\begin{itemize}
  \item \textbf{Value} sorts stocks on book-to-market\index{book-to-market ratio} once a year. Book-to-market ratios change slowly, so most stocks stay on the same side from one year to the next.
  \item \textbf{Momentum}\index{momentum|textbf} sorts stocks each month on their return from 12 months to 2 months ago. Each month's window shares ten of its eleven months with the previous one. Even so, enough stocks cross the \textbf{breakpoints}, the cut-offs between groups, to replace about a third of the portfolio each month.
  \item \textbf{Short-term reversal}\index{short-term reversal|textbf} sorts stocks each month on the previous month's return. Last month's return has almost no persistence\index{persistence}, so nearly the whole portfolio is replaced every month.
\end{itemize}

\textbf{From gross to net.} A \textbf{gross} return is the return before costs, and a \textbf{net} return is the return after them. Consider a long--short portfolio with one dollar on each side and one-sided turnover $T$. Each month it sells $T$ and buys $T$ on each side, so it trades $4T$ dollars in total. At a one-way cost $c$, its monthly cost and net return are
\begin{equation}
  \text{cost} = 4\,T\,c, \qquad R^{\text{net}}_t = R^{\text{gross}}_t - 4\,T\,c - \frac{\text{borrow fee}}{12}
  \label{eq:04-net}
\end{equation}
where the borrowing fee is quoted per year.

The \textbf{break-even cost}\index{transaction costs!break-even|textbf} is the one-way cost at which the average net return falls to zero:
\begin{equation}
  c^* = \frac{\text{average gross return} - \text{other costs}}{4T}
  \label{eq:04-breakeven}
\end{equation}
It answers the most useful question about a strategy's costs: how cheaply must it trade to be worth trading? For example, a strategy with a turnover of 10\% a month that earns 0.4\% a month gross can pay up to $0.4/(4 \times 0.10) = 1\%$ per trade before its profit is gone.

Novy-Marx and Velikov (2016) estimated turnover and costs for value-weighted versions of 23 anomalies. The table below reports their full-sample figures for the three strategies in this chapter, from Table 3 of the working-paper version. The break-even costs are computed from their gross returns with Equation~\eqref{eq:04-breakeven}.

\begin{center}
\small
\begin{tabularx}{0.9\linewidth}{@{}>{\raggedright\arraybackslash}X>{\raggedright\arraybackslash}p{0.13\linewidth}>{\raggedright\arraybackslash}p{0.13\linewidth}>{\raggedright\arraybackslash}p{0.13\linewidth}>{\raggedright\arraybackslash}p{0.13\linewidth}>{\raggedright\arraybackslash}p{0.12\linewidth}@{}}
\toprule
Strategy & One-sided monthly turnover & Gross return (\% per month) & Estimated cost (\% per month) & Net return (\% per month) & Break-even one-way cost \\
\midrule
Value (book-to-market) & 2.9\% & 0.47 & 0.05 & 0.42 & about 4\% \\
Momentum & 34.5\% & 1.33 & 0.65 & 0.68 & about 1\% \\
Short-term reversal & 90.9\% & 0.37 & 1.65 & $-1.28$ & about 0.1\% \\
\bottomrule
\end{tabularx}
\end{center}

The three strategies sit on very different sides of realistic costs (Figure~\ref{fig:04-breakeven}). Value could pay several percent per trade and still make a profit. Momentum survives costs of up to about one percent per trade, which a careful institution can achieve but a careless one may not. Short-term reversal is profitable on paper, but it would need to trade at about ten basis points per trade. That is below what most investors pay for the small, recently volatile stocks it buys and sells. Its paper profit is real, but it is tiny compared with the trading it requires.

\begin{figure}[htbp]
  \centering
  \begin{tikzpicture}[
    % Gross return (% a month) and one-sided monthly turnover, from the table.
    % Net return = gross - 4 * turnover * cost; break-even = gross / (4 * turnover).
    declare function={
      valG = 0.47;  valT = 0.029;
      momG = 1.33;  momT = 0.345;
      revG = 0.37;  revT = 0.909;
    }]
  \begin{axis}[
    width=0.61\linewidth, height=0.38\linewidth,
    axis lines=left,
    axis line style={-stealth, thin},
    xmin=0, xmax=1.3, ymin=-1.6, ymax=1.45,
    xtick={0,0.2,0.4,0.6,0.8,1,1.2},
    ytick={-1,0,1},
    xticklabel style={/pgf/number format/fixed},
    xlabel={One-way trading cost (\%)},
    ylabel={Net return (\% a month)},
    clip=false,
  ]
    \addplot[gray!60, thin, domain=0:1.3, samples=2] {0};

    % Value
    \addplot[fenavy, very thick, domain=0:1.3, samples=2] {valG - 4*valT*x};
    % Momentum
    \addplot[thick, dashed, domain=0:1.3, samples=2] {momG - 4*momT*x};
    % Short-term reversal (stops at the bottom of the chart)
    \addplot[very thick, densely dotted, domain=0:0.54, samples=2] {revG - 4*revT*x};

    % Break-even points
    \fill (axis cs:{momG/(4*momT)},0) circle (2.2pt);
    \fill (axis cs:{revG/(4*revT)},0) circle (2.2pt);
    \node[anchor=north east, font=\small] at (axis cs:{momG/(4*momT)-0.02},-0.03)
      {break-even $\approx 1\%$};
    \node[anchor=south west, font=\small] at (axis cs:{revG/(4*revT)+0.03},0.03)
      {$\approx 0.1\%$};

    % Line labels
    \node[fenavy, anchor=west, font=\small] at (axis cs:1.31,{valG-4*valT*1.3+0.05}) {value};
    \node[anchor=west, font=\small] at (axis cs:1.31,{momG-4*momT*1.3}) {momentum};
    \node[anchor=west, font=\small] at (axis cs:0.48,-1.05) {short-term reversal};
  \end{axis}
  \end{tikzpicture}
  \caption{Net return against one-way trading cost, from the turnover and gross
    returns in the table. Each line falls at a rate of four times the strategy's
    turnover and crosses zero at the break-even cost. Value's line is nearly
    flat and crosses zero off the chart, at about 4\%.}
  \label{fig:04-breakeven}
\end{figure}

A ranking of strategies by their gross Sharpe ratio\index{Sharpe ratio} or $t$-statistic\index{t-statistic} therefore says little about which one an investor should prefer. The net ranking depends on turnover as much as on the gross premium.

\section{Market impact and capacity}
\label{sec:04-capacity}

A spread costs roughly the same amount for every dollar traded. Market impact does not. The more a strategy trades relative to the market's volume, the more each dollar costs. This is what gives every strategy a \textbf{capacity}\index{capacity|textbf}: a size beyond which its net return shrinks toward zero.

\textbf{Why prices move.} Kyle (1985) modeled market makers who cannot tell informed orders from uninformed ones. To protect themselves, they move the price in the direction of every order. Large orders are more likely to be informed and harder to absorb, so they move the price more. Studies of large institutional orders find that impact grows less than in proportion to the order's size. Almgren, Thum, Hauptmann and Li (2005) estimated that it grows with the order's size raised to a power of about 3/5, and T\'oth et al.\ (2011) found a power close to 1/2 across many markets. The second result, the \textbf{square-root law}, is the common working model:
\begin{equation}
  \text{impact per dollar} \approx Y\,\sigma_{daily}\,\sqrt{\frac{Q}{V}}
  \label{eq:04-impact}
\end{equation}
Here $Q$ is the order size, $V$ is the day's trading volume, $\sigma_{daily}$ is the stock's daily volatility and $Y$ is a constant close to one. The ratio $Q/V$ is the \textbf{participation rate}\index{participation rate|textbf}, the order's share of the day's trading. Trading 1\% of a day's volume in a stock with 2\% daily volatility costs about $2\% \times \sqrt{0.01} = 0.2\%$ in impact. Trading 10\% costs about $2\% \times \sqrt{0.10} \approx 0.6\%$ (Figure~\ref{fig:04-impact}).

\begin{figure}[htbp]
  \centering
  \begin{tikzpicture}[
    % Impact (%) = Y * sigma * sqrt(Q/V), with Y = 1 and sigma = 2%.
    % x is Q/V in percent, so impact = 2 * sqrt(x/100) = 0.2 * sqrt(x).
    declare function={impact(\x) = 0.2*sqrt(\x);}
  ]
  \begin{axis}[
    width=0.55\linewidth, height=0.35\linewidth,
    axis lines=left,
    axis line style={-stealth, thin},
    xmin=0, xmax=20.5, ymin=0, ymax=1,
    xtick={0,2,...,20},
    ytick={0,0.2,0.4,0.6,0.8,1},
    yticklabel style={/pgf/number format/fixed},
    tick label style={font=\small},
    xlabel={Order size as a share of daily volume, $Q/V$ (\%)},
    ylabel={Impact per dollar traded (\%)},
    samples=200, domain=0:20,
    clip=false,
  ]
    % If impact were linear (for comparison)
    \addplot[gray!70, densely dashed, thin, samples=2] {0.0215*x};
    \node[gray!40!black, font=\small, anchor=west] at (axis cs:20.2,0.43)
      {if impact were linear};

    % Square-root impact
    \addplot[fenavy, very thick] {impact(x)};
    \node[fenavy, font=\small, anchor=west] at (axis cs:20.2,{impact(20)})
      {square root};

    % Marked points
    \fill[fenavy] (axis cs:1,0.2) circle (2.2pt);
    \draw[gray!50] (axis cs:1.2,0.21) -- (axis cs:4.6,0.34);
    \node[font=\small, anchor=west] at (axis cs:4.6,0.34) {1\% of volume: 0.2\%};

    \fill[fenavy] (axis cs:10,0.632) circle (2.2pt);
    \node[font=\small, anchor=north west] at (axis cs:10.3,0.6)
      {10\% of volume: 0.6\%};
  \end{axis}
  \end{tikzpicture}
  \caption{Square-root market impact for a stock with 2\% daily volatility,
    from Equation~\eqref{eq:04-impact} with $Y = 1$. Ten times the order size
    costs about three times as much per dollar, not ten times. Total
    impact cost, per dollar times dollars, still rises faster than the order.}
  \label{fig:04-impact}
\end{figure}

\textbf{From impact to capacity.} A strategy that manages assets $A$ with turnover $T$ trades about $2TA$ on each side every month, spread across its holdings. As $A$ grows, the participation rate $Q/V$ rises, impact per dollar rises with its square root, and total impact cost rises faster than the assets. Net return therefore falls as the fund grows, slowly at first and then quickly. Capacity can be defined as the size at which the net return is halved, or the size at which it reaches zero. An investor would want to stop well before either.

Three features of a strategy set its capacity.

\begin{enumerate}
  \item \textbf{Turnover.}\index{turnover} Every dollar that is not traded is a dollar that does not move prices.
  \item \textbf{The liquidity of its holdings.} Strategies that concentrate in small stocks reach their limits first.
  \item \textbf{The gross premium.} A larger premium can absorb more cost before it is gone.
\end{enumerate}

Capacity also has an industry-wide side. If many managers follow similar signals, their combined trades determine the impact, and the strategy's capacity is shared among everyone who trades it. Berk and Green (2004) argued that in a competitive market, money flows into skilled managers until their returns after costs are no better than those of the alternatives. Rising costs at larger sizes are one of the forces that bring this about.

The capacity estimates in this chapter's lab rest on stated assumptions about volume, spreads and the impact constant. They show how capacity arises and what it depends on. They are not a forecast for any real fund.

\section{Designing for costs}
\label{sec:04-design}

A paper portfolio is built as if trading were free, so it trades whenever the signal moves, however slightly. A portfolio designed for costs trades only when the expected gain is larger than the cost. A few simple rules capture most of the benefit.

\begin{itemize}
  \item \textbf{Buy/hold spreads, or banding.} Buy a stock only when it enters the top 30\% on the signal, but keep holding it until it falls out of, say, the top 50\%. Stocks that wobble around the breakpoint then no longer generate trades. Novy-Marx and Velikov (2016) found this to be the single most effective simple way to reduce costs.
  \item \textbf{Slower rebalancing\index{rebalancing}.} Rebuilding the portfolio once a quarter rather than once a month cuts turnover for slow signals at little loss of signal. For fast signals it gives up most of the premium.
  \item \textbf{Partial trading.} G\^arleanu and Pedersen (2013) studied an investor whose returns are partly predictable\index{return predictability} and whose trading costs rise with the square of the amount traded, called \textbf{quadratic trading costs}\index{transaction costs!quadratic|textbf}. They showed that the best policy trades only part of the way toward the target each period, and aims at a target that anticipates where the signal is heading.
  \item \textbf{Netting and combining signals.} A portfolio that uses several signals can cross trades internally. A stock that momentum wants to sell and value wants to buy need not be traded at all.
  \item \textbf{Liquidity screens.}\index{liquidity|textbf} Restricting the portfolio to larger, more liquid stocks raises capacity and lowers costs. It usually gives up some of the gross premium, because many anomalies are strongest in small stocks.
\end{itemize}

Each rule gives up a little signal in exchange for much less trading. Whether the exchange pays depends on how fast the signal decays. Banding costs little for momentum, whose ranks change gradually. It cannot rescue short-term reversal, whose signal is gone within a month.

\textbf{Cost mitigation is also a search.} Every design choice, such as a band width, a rebalancing frequency or a liquidity screen, can be tuned to make a backtest look better. The warnings about selection and overfitting\index{overfitting} in Section~\ref{sec:01-backtests} therefore apply in full. Chen and Velikov (2023) examined 204 anomalies and accounted for effective spreads, for the decline in returns after publication\index{anomaly!post-publication decay}, and for the modern era of electronic trading. They found that the average anomaly's expected return, net of these effects, was only about 4 basis points a month, and this was before counting price impact. Strategies that showed impressive net returns in sample after cost mitigation did not escape that conclusion. A cost-aware design should be fixed in advance and judged out of sample, like every other part of a strategy.

\begin{callout}
\textbf{What Failed: buying last month's losers.} Jegadeesh (1990) and Lehmann (1990) found that stocks with low returns over the past week or month tended to rebound, and that stocks with high returns tended to fall back. The pattern is strong and statistically robust on paper, and the Fama--French\index{Fama--French factors} short-term reversal factor still has a positive average return. But the strategy replaces almost its whole portfolio every month. It also concentrates in stocks that have just moved sharply, which are often small and illiquid, exactly where spreads are widest. Novy-Marx and Velikov (2016) estimated that its costs exceeded its gross return by more than a percentage point a month. Nagel (2012) showed that much of the reversal return can be understood as payment for supplying liquidity to impatient traders. That payment goes to market makers who trade at low cost, not to investors who must pay the spread. As a strategy for an ordinary investor, the claim that buying recent losers is profitable \textbf{fails} once costs are counted.
\end{callout}

\section{Lab~\thechapter: What survives the trip from paper to portfolio?}
\label{sec:04-lab}

\begin{callout}
\textbf{Lab~\ref{ch:04} at a glance.} \emph{Claim on trial:} a strategy that is profitable in a backtest is profitable to trade. \emph{Data:} three long--short factors from the Kenneth R. French Data Library, value (HML), momentum (Mom) and short-term reversal (ST Rev), from July 1963, plus a simulated stock universe. \emph{Run time:} a few seconds. \emph{Code:} repository folder \texttt{ch05-costs/}, with a guided notebook.
\end{callout}

The lab takes the three factors and asks what each would have earned after costs.

\textbf{Part A: gross performance.} The lab reports average returns, $t$-statistics and Sharpe ratios before costs, as the factors are published. On paper, all three look like discoveries.

\textbf{Part B: turnover.} The French library does not publish holdings, so turnover comes from two sources. The first is Novy-Marx and Velikov's estimates for value-weighted versions of the same strategies. The second is a simulated universe of 1,000 stocks over 40 years. Its returns have fat tails\index{fat tails}, its companies have a realistic spread of sizes, and its portfolios are built exactly as the Fama--French factors are: value-weighted 30/70 sorts, with value updated each June and momentum and reversal each month. Because the simulated universe has no delistings\index{delisting} or new listings, its turnover is a lower bound. In the version shipped with the lab, it gives about 2\% a month for value, 25\% for momentum and 76\% for reversal. These are in the same order and of similar size as the published estimates, which shows that turnover follows mechanically from how each signal is built.

\textbf{Part C: costs.} The lab computes net returns at one-way costs from 0 to 100 basis points, with a borrowing fee of 0.25\% a year. It reports each strategy's break-even cost, and it ranks the three strategies by Sharpe ratio before and after costs at a headline cost of 20 basis points. Lab Figure~\ref{ch:04}.1 plots each strategy's net return against cost, with its break-even point marked. Lab Figure~\ref{ch:04}.2 shows the growth of a dollar on paper and after costs.

\textbf{Part D: capacity.} The simulated universe is scaled to a \$50 trillion market. Each company's daily volume is 0.4\% of its market value, half-spreads run from 2 basis points for the largest stocks to 50 for the smallest, and impact follows the square-root law. The lab traces each strategy's net return as its assets grow from \$10 million to \$1 trillion. The assumptions are stated and can be changed, and the exercises ask which of them matters most.

\textbf{Part E: cost mitigation, with known truth.} In the simulated universe the expected returns are known, so cost-mitigation rules can be compared without sampling noise. The lab compares the standard sort with buy/hold spreads of increasing width. For momentum at a one-way cost of 50 basis points, holding stocks until they leave the top 50\% cuts turnover roughly in half, from about 25\% to 12\% a month. It gives up about a fifth of the expected gross return, and it raises the expected net return from about 0.30\% to 0.41\% a month. For reversal, no band width rescues the strategy at that cost, because the signal decays too fast for patience to pay.

The lab, its tests and a guided notebook are in \texttt{ch05-costs} in the book's repository. The README records what has been tested and how. It notes that the parsers for the momentum and reversal files have been tested on the library's layout but not yet on the live files.

\section{Verdict}
\label{sec:04-verdict}

The verdict on the claim is that it \textbf{fails} as a general rule. Whether a profitable backtest is profitable to trade depends on how much the strategy trades, what it trades, and how much money is behind it. Slow strategies in liquid stocks survive realistic costs comfortably. Fast strategies in small stocks often do not survive them at all. The claim holds only for strategies whose break-even costs sit well above what their investors actually pay, at the size they intend to run.

\section{What investors can use}
\label{sec:04-investors}

Five questions turn the chapter into a checklist for any strategy an investor is shown.

\begin{enumerate}
  \item \textbf{What is the turnover?} Ask for one-sided monthly turnover on both sides. If the answer is not known, the strategy has not been evaluated as an investment.
  \item \textbf{What is the break-even cost, and what will you pay?} Compare the break-even one-way cost with realistic costs for the stocks the strategy holds, in the period it will be traded. A strategy that needs costs below ten basis points is a bet on execution, not on the signal.
  \item \textbf{What does the short side cost?} Ask how much of the return comes from the short side, what borrowing fees were assumed, and how the strategy would behave if hard-to-borrow stocks could not be shorted.
  \item \textbf{At what size?} Ask for net returns at the assets the manager plans to run, and how market impact was estimated. Capacity claims without stated assumptions are marketing.
  \item \textbf{Were cost-saving rules chosen in advance?} Bands, rebalancing schedules and liquidity screens are sensible, but each is a parameter. Ask whether they were fixed before the backtest, and how the strategy performed after they were fixed.
\end{enumerate}

The broader lesson is that costs change the ranking of strategies, not just their level. A modest, slow premium that can be harvested cheaply may be worth more than a spectacular, fast one that cannot. Chapters~\ref{ch:05} to~\ref{ch:07} test the best-known anomalies, and every verdict there is reported after costs as well as before.

\section*{Exercises}

\subsection*{Conceptual (all readers)}

\begin{enumerate}
  \item A long--short strategy has one-sided monthly turnover of 20\% and earns 0.6\% a month gross. What is its break-even one-way cost? What is its net return at a one-way cost of 30 basis points?
  \item Explain why a value-weighted portfolio whose membership never changes has zero turnover, while an equal-weighted one with the same membership does not.
  \item Why does market impact give a strategy a capacity, while a fixed spread does not?
  \item A strategy earns most of its paper return on the short side. List three reasons its implementable return could be much lower than its paper return.
\end{enumerate}

\subsection*{Lab extensions (technical readers)}

\begin{enumerate}
  \setcounter{enumi}{4}
  \item Re-run the lab from May 2001, after the move to decimal prices, with a one-way cost of 10 basis points. Does short-term reversal survive in the low-cost era?
  \item Re-run Part C with a borrowing fee of 2\% a year. Which strategy is most exposed to shorting costs, and why?
  \item In Part D, halve and double the impact constant and the number of days over which trades are spread. Which assumption moves capacity most? What evidence would you want before trusting it?
  \item In Part E, add a rule that rebalances momentum only once a quarter. Compare its turnover and expected net return with banding.
  \item Implement the Corwin--Schultz high--low spread estimator and apply it to daily data for a few stocks you can download. How does it compare with their quoted spreads?
  \item Choose an anomaly from Chapter~\ref{ch:07}'s list. Estimate its turnover from its construction, compute its break-even cost, and judge whether it is likely to survive costs.
\end{enumerate}
